% Cote 119: offprints, reports, courses, CV and letters (typed and printed part): archive cote 119, site page 69 (PDF page 70).
% offprint. Draft compared: none.
% AI-transcribed by Claude Opus 5.5 (claude-opus-5-5, adaptive thinking, effort high), Message Batches API, batch r119_full (msgbatch_01DwZNm5Y7YsKBWyAN4sFs8w), 2026-11-13; not yet checked by eye.
\scanpage{119/69}
\mspage{C. R. Acad. Sc. Paris, t. 272, p. 258-261 (18 janvier 1971). \quad Série A}

GÉOMÉTRIE ALGÉBRIQUE. --- Représentations linéaires de schémas
en groupes affines. Note (*) de M. Neantro Saavedra Rivano, transmise
par M. Henri Cartan.

La correspondance qui, à un schéma en groupes affine, associe la catégorie de
ses représentations linéaires, permet d'établir un dictionnaire entre groupes affines
et certaines $\otimes$-catégories.

1. La catégorie Rep(G).
1.1. Soient S un schéma, G un S-monoïde affine [($^1$), I, 4]. On note
Rep(G) la catégorie des représentations linéaires de G dans des $\mathcal{O}_S$-Modules
quasi-cohérents (i.e., G-$\mathcal{O}_S$-Modules quasi-cohérents à gauche),
\[ \omega^G : \mathrm{Rep}(G) \to \mathrm{Qcoh}(S) \]
le foncteur « oubli de l'action G ». On notera que Rep(G) est canoni-%
quement une $\otimes$-catégorie ACU, $\omega^G$ un $\otimes$-foncteur ACU ($^2$), et que la
formation de la $\otimes$-catégorie Rep(G) au-dessus de Qcoh(S) est fonc-%
torielle en le monoïde G. Si $u : G \to G'$ est un morphisme de S-monoïdes
affines, on notera $\mathrm{Rep}(u) : \mathrm{Rep}(G') \to \mathrm{Rep}(G)$ le $\otimes$-foncteur ACU qu'il
induit. \margin{\unsure{/}}

Supposons que G soit S-plat, ce qui pour S affine équivaut à la commu-%
tation de $\omega^G$ avec les limites projectives finies, et entraîne que Rep(G)
est une catégorie abélienne. On appelle foncteur fibre sur Rep(G) un
$\otimes$-foncteur ACU$\Gamma(S, \mathcal{O}_S)$-linéaire $\omega : \mathrm{Rep}(G) \to \mathrm{Qcoh}(S)$ qui soit fidèle,
exact et qui commute avec les limites inductives. En prenant comme
morphismes les $\otimes$-morphismes unifères, on obtient une catégorie Fib(G)
dont l'ensemble des classes d'isomorphisme d'objets sera noté fib(G).
1.2. Soient $C$ une $\otimes$-catégorie ACU, $a$, $b$ des $\otimes$-foncteurs
ACU $C \to \mathrm{Qcoh}(S)$. On note $Hom^{\otimes}(a, b)$ le préfaisceau sur la catégorie
$\mathrm{Sch}_{/S}$ qui à un S-schéma S$'$ associe l'ensemble
\[ Hom^{\otimes}(a, b) = \mathrm{Hom}^{\otimes}(a_{S'}, b_{S'}) \]
des $\otimes$-morphismes unifères de $a_{S'}$ dans $b_{S'}$, la notation $a_{S'}$ désignant le
$\otimes$ - foncteur ACU composé de $a$ et de l'extension des scalaires
$\mathrm{Qcoh}(S) \to \mathrm{Qcoh}(S')$. On définit de même $Isom^{\otimes}(a, b)$, $End^{\otimes}(a)$, $Aut^{\otimes}(a)$.
Il s'agit en fait de faisceaux pour la topologie fidèlement plate quasi
compacte (fpqc) [($^1$), VII, 4].

Revenant à la situation de 1.1, la définition même de G-$\mathcal{O}_S$-Module
montre qu'on dispose d'un morphisme canonique de faisceaux en monoïdes
sur S :
\[ (1.2.1) \qquad G \to End^{\otimes}(\omega^G) \]

\mspage{\hand{69}}
